\frontmatter

% ============================================================
% Halaman Sampul
% ============================================================
\begin{titlepage}
  \centering
  \vspace*{1cm}
  {\LARGE\bfseries Logika Matematika\par}
  \vspace{0.5cm}
  {\large Buku Ajar Berbasis Outcome-Based Education (OBE)\par}
  \vspace{1cm}
  {\large Mata Kuliah: Logika Matematika\par}
  {\large Kode: IT-231\par}
  \vspace{2cm}
  {\large Program Studi S1 Teknik Informatika\par}
  {\large Fakultas Teknologi Informasi\par}
  {\large Universitas Bale Bandung\par}
  \vfill
  {\large 2026\par}
\end{titlepage}

\cleardoublepage

% ============================================================
% Kata Pengantar
% ============================================================
\chapter*{Kata Pengantar}
\addcontentsline{toc}{chapter}{Kata Pengantar}

Puji syukur kehadirat Tuhan Yang Maha Esa atas terselesaikannya buku ajar \textit{Logika Matematika} ini. Buku ini disusun dengan pendekatan \textbf{Outcome-Based Education (OBE)}, yang berfokus pada pencapaian kompetensi terukur sesuai dengan Capaian Pembelajaran Lulusan (CPL) dan Capaian Pembelajaran Mata Kuliah (CPMK) Program Studi Teknik Informatika Universitas Bale Bandung.

Logika Matematika merupakan mata kuliah fundasional yang mengajarkan pola pikir deduktif yang terstruktur dan sangat diperlukan sebagai basis analisis algoritma maupun komputasi. Paradigma dalam buku ini disusun agar mahasiswa mampu mengaplikasikan landasan logika proposisi, logika predikat, hingga aljabar Boolean dalam menyelesaikan perancangan komputasi.

Buku ini dirancang untuk mendukung pembelajaran dengan pendekatan student-centered learning. Setiap bab dilengkapi dengan Sub-CPMK terukur, penjelasan materi, aktivitas dan latihan pemecahan masalah, instrumen asesmen diri, dan disajikan memuat 17 bab untuk 16 kali pertemuan plus soal evaluasi komprehensif.

Kami berharap buku ini dapat menjadi instrumen belajar yang efektif.

\vspace{1cm}
\begin{flushright}
Penyusun
\end{flushright}

\cleardoublepage

% ============================================================
% Cara Menggunakan Buku Ini
% ============================================================
\chapter*{Cara Menggunakan Buku Ini}
\addcontentsline{toc}{chapter}{Cara Menggunakan Buku Ini}

Buku ajar ini dirancang dengan pendekatan OBE untuk memaksimalkan pencapaian pembelajaran Anda. 

\section*{Struktur Buku}

\begin{itemize}
\item \textbf{Bab 1: Pendahuluan dan Orientasi} Memperkenalkan kerangka OBE dan relevansi Logika Matematika di tingkat sarjana.
\item \textbf{Bab 2: Pengenalan mengenai Logika Matematika} Menyajikan sejarah, landasan filosofi, serta konsep formal matematika secara luas sebelum masuk ke teknis.
\item \textbf{Bab 3 -- Bab 16: Unit Materi Inti} Memuat seluruh bahasan matakuliah (Logika Proposisi, Predikat, Boolean, dan Induksi). Bab disusun selaras dengan rencana mingguan RPS.
\item \textbf{Bab 17: Soal Evaluasi} Menyajikan kumpulan instrumen Ujian Tengah dan Akhir Semester.
\end{itemize}

\section*{Komponen dalam Setiap Bab}

\begin{enumerate}
  \item \textbf{Sub-CPMK}: Kompetensi yang harus Anda kuasai pada bab terpilih.
  \item \textbf{Materi Pokok}: Rincian akademis dan contoh kasus komputasi.
  \item \textbf{Aktivitas dan Latihan}: Praktik mereduksi logika, penyederhanaan tabel kebenaran atau penarikan simpulan matematis.
  \item \textbf{Checklist}: Instrumen mandiri untuk meyakinkan kompetensi.
\end{enumerate}

\cleardoublepage

% ============================================================
% Identitas Mata Kuliah
% ============================================================
\chapter*{Identitas Mata Kuliah}
\addcontentsline{toc}{chapter}{Identitas Mata Kuliah}

\begin{tabular}{ll}
  Program Studi & : S1 Teknik Informatika \\
  Nama Mata Kuliah & : Logika Matematika \\
  Kode Mata Kuliah & : IT-231 \\
  Bobot SKS & : 2 SKS \\
\end{tabular}

\vspace{1cm}

\section*{Capaian Pembelajaran Lulusan (CPL)}

\begin{enumerate}
  \item \textbf{CPL01:} Kemampuan menganalisis persoalan computing yang kompleks serta menerapkan prinsip-prinsip computing dan disiplin ilmu relevan lainnya untuk mengidentifikasi solusi, dengan mempertimbangkan wawasan perkembangan ilmu transdisiplin.
  \item \textbf{CPL02:} Memiliki pengetahuan yang memadai terkait dengan cara kerja sistem komputer dan mampu merancang dan mengembangkan berbagai algoritma metode untuk memecahkan masalah.
\end{enumerate}

\section*{Capaian Pembelajaran Mata Kuliah (CPMK)}

\begin{enumerate}
  \item \textbf{CPMK1:} Mahasiswa mampu menganalisis validitas argumen dan persoalan komputasi dasar dengan menerapkan prinsip-prinsip logika proposisional, tabel kebenaran, dan penarikan kesimpulan.
  \item \textbf{CPMK2:} Mahasiswa mampu memodelkan persoalan yang lebih kompleks ke dalam bentuk ekspresi logika predikat.
  \item \textbf{CPMK3:} Mahasiswa mampu menerapkan hukum-hukum aljabar Boolean untuk menyederhanakan fungsi logika dan merepresentasikannya ke dalam bentuk gerbang logika dasar.
  \item \textbf{CPMK4:} Mahasiswa mampu merancang argumen matematis dan membuktikan kebenaran suatu metode menggunakan teknik pembuktian formal dan induksi matematika.
\end{enumerate}

\cleardoublepage

% ============================================================
% Daftar Isi
% ============================================================
\phantomsection
\addcontentsline{toc}{chapter}{Daftar Isi}
\tableofcontents
\cleardoublepage
